\section{Ordinary modules and integral auxiliary control}
\label{sec:sr-ordinary-control}
We construct the ordinary function modules, their coefficient change and
classical specialization maps, and their integral old/new operators.
These supply the module-theoretic part of common support patching.
The attached Galois representations, their local compatibility, the
arithmetic Selmer bounds and the seed remain separate requirements.
The same constructions apply to ranks three and four.


\subsection{Construction of the ordinary modules and their finite-level control}
\label{sec:ordinary-module-construction}

This subsection constructs the finite-projective modules used by the
auxiliary-level argument.  It does not take \srcid{LevelAlt} or
\srcid{LevMod} as a premise.  The arithmetic setting is the definite
unitary setting of \srcid{HidaSetting}: \(L\) is a CM field,
\(G/L^+\) is the unitary group of the standard positive definite
Hermitian form of rank \(n\geq1\), and all places used for the ordinary
operators and for varying the auxiliary level are split places.
At the distinguished fixed auxiliary place the level is contained in
the first principal congruence subgroup at an odd rational prime \(p\),
with \(p\ne\ell\) and \(p\) unramified in \(L\).  These are exactly the
smallness conditions used below.  The construction is valid in every
rank; in the generic application \(n=3\).

Let \(A\) be a complete Noetherian local coefficient algebra with finite
residue field and maximal ideal \(\mathfrak m\).  Let \(A^m\) be the
coefficient lattice, with the integral tame action denoted by \(\tau\).
Write \(G_{\mathrm a}\) for the factors away from \(\ell\), \(G_\ell\)
for the product of the split factors above \(\ell\), and
\(\Gamma=G(L^+)\) for the rational subgroup.  On \(G_\ell\) fix the
Iwahori data
\[
 I(c)=N\,T\,N^-_c,\qquad c\geq0,
\]
with their usual opposite-unipotent filtration \(N^-_c\), diagonal
congruence subgroups \(T_c\), and a strictly positive central element
\(z\) in the positive diagonal monoid.  Products and factorizations here
are interpreted in the Iwahori decomposition; in particular
\(I(c)\subset I(0)\), \([I(0):I(c)]<\infty\), and multiplication by \(z\)
contracts the opposite-unipotent filtration by at least one step.
The character \(\theta:T\to A^\times\) is continuous in the precise
sense
\[
 \forall k\ \exists c\quad
 \theta(t)-1\in\mathfrak m^k\quad(t\in T_c).
 \tag{OC1}
\]
The universal weight characters occurring in the application have this
property.

For a tame compact open subgroup \(U\), let \(S(U)\) be the \(A\)-module
of functions
\[
 f:G_{\mathrm a}\times G_\ell\longrightarrow A^m
\]
which are left \(\Gamma\)-invariant, satisfy the tame
\(\tau\)-equivariance, are right \(N\)-invariant and right \(T\)-equivariant
with character \(\theta^{-1}\), and are uniformly continuous along
\(N^-_c\) modulo every \(\mathfrak m^k\).
The inverse in \(\theta^{-1}\) is the convention in the function model.
Put
\[
 V_k=\{f\in S(U):f(g)\in(\mathfrak m^k)^m\ \hbox{for every }g\},
 \qquad T_z=U_z.
\]
Here \(U_z\) is the integral double-coset sum, with the normalization of
the source Iwahori data.  It preserves \(S(U)\) and every \(V_k\).
Define the ordinary module by
\[
 O(U)=
 \{f\in S(U):\ \forall k\ \exists r\geq1,\
                 T_z^r f-f\in V_k\}.
 \tag{OC2}
\]

\begin{srlemma}[Smallness of the actual arithmetic levels]
\label{lem:ordinary-neat-level}
For every collection of compact open subgroups at the varying
auxiliary places, the tame level in the preceding construction is
neat in the following explicit sense:
\[
 \gamma g(u,b)=g,\quad
 \gamma\in\Gamma,\ u\in U,\ b\in I(0)
 \quad\Longrightarrow\quad u=b=1.
 \tag{OC3}
\]
Consequently arbitrary values in \(A^m\) can be prescribed on distinct
double cosets for \(\Gamma\backslash
(G_{\mathrm a}\times G_\ell)/(U\times I(c))\).
\end{srlemma}

\begin{proof}
Represent \(\gamma\) by a matrix \(M\in\mathrm{GL}_n(L)\).
At every archimedean embedding the matrices \(M^r\) are unitary.
At the varying auxiliary places, compactness implies that some positive
power of the local matrix belongs to a conjugate of
\(\mathrm{GL}_n(\mathcal O_{L,w})\): its orbit on the discrete set of
lattices is finite.  Taking a common multiple for the finitely many
varying places, and using the defining integral levels at the remaining
places, we find a positive integer \(a\) such that \(M^a\) is locally
conjugate to an integral invertible matrix at a selected place above
each finite place of \(L^+\).

Thus the coefficients of the characteristic polynomial of \(M^a\)
are integral at those selected places, as are their conjugates:
unitarity gives \(\overline{M^a}{}^t=(M^a)^{-1}\), and the latter is
also locally conjugate to an integral invertible matrix.
For a coefficient \(x\), both
\(x+\overline x\) and \(x\overline x\) belong to \(\mathcal O_{L^+}\);
the equation
\(X^2-(x+\overline x)X+x\overline x=0\) proves that \(x\) is integral
over \(\mathbb Z\).  Cayley--Hamilton now puts all powers of \(M^a\)
in a finitely generated \(\mathcal O_L\)-module of matrices.
After one fixed denominator is cleared, every entry of every power
is an algebraic integer bounded at all archimedean embeddings.
The Minkowski lattice has only finitely many such elements.
The set of powers of \(M^a\), and hence of \(M\), is finite.

At the distinguished auxiliary place the matrix is conjugate to a
finite-order element of \(1+pM_n(\mathcal O_{L,w})\).
This subgroup is torsion free when \(p\) is odd and unramified.
Indeed, if \(1+C\ne1\) and the minimum valuation of its entries is
\(r\geq1\), raising to a power prime to \(p\) retains that minimum,
whereas raising to the \(p\)-th power gives minimum \(r+1\).
For the latter assertion the first term \(pC\) has minimum \(r+1\);
all remaining binomial terms have larger minimum, since \(p\geq3\).
No positive power can therefore equal \(1\).
It follows that \(M=1\); the stabilizer identity then gives \(u=b=1\).

The finite double-coset assertion needed here follows from the definite
unitary finite-adelic orbit argument already proved in the arithmetic
finiteness subsection.  Passing from \(I(0)\) to \(I(c)\) preserves
finiteness because its index is finite.  In view of (OC3), the values on
the resulting double-coset representatives have no stabilizer
compatibility condition.
\end{proof}

\begin{srlemma}[Ordinary projector and finite freeness]
\label{lem:ordinary-projector-free}
For every arithmetic level above there is an \(A\)-linear projection
\[
 e_U:S(U)\longrightarrow O(U),\qquad
 e_U f=\lim_{r\to\infty}T_z^{\,r!}f
 \tag{OC4}
\]
in the valuewise \(\mathfrak m\)-adic filtration.  The module \(O(U)\)
is finite free over \(A\), and
\[
 O(U)\cap V_k=\mathfrak m^kO(U).
 \tag{OC5}
\]
Every positive diagonal \(U\)-operator is invertible on \(O(U)\).
All assertions hold uniformly for arbitrary compact open choices at
the finitely many varying auxiliary places; the finite free rank may
depend on the level.
\end{srlemma}

\begin{proof}
We give the finite approximation and algebra arguments, since this
avoids importing an unspecified ordinary control theorem.

Fix \(k\), choose \(c\) as in (OC1), and choose representatives
\(x_1,\ldots,x_b\) for the finite double-coset set with \(I(c)\).
Evaluation gives
\[
 E_k:S(U)\longrightarrow A^{bm}.
\]
There is an \(A\)-linear interpolation
\(s_k:A^{bm}\to S(U)\) with \(E_ks_k=1\), and its values are
\(N^-_c\)-invariant modulo \(\mathfrak m^k\).
Here is the explicit construction.  For \(b\in I(0)\), let \(t(b)\)
be the diagonal factor in \(b=y\,t(b)\,x\), with
\(y\in N^-_0\), \(x\in N\).  Its uniqueness gives
\[
 t(sb)=s\,t(b),\quad t(bs)=t(b)s,\quad t(bx)=t(b).
\]
The Iwahori interchange relation further gives
\(t(by)=t(b)s\) for some \(s\in T_c\) whenever \(y\in N^-_c\).
For prescribed vectors \(v_i\), set
\[
 s_k(v)\bigl(\gamma x_i(u,b)\bigr)
       =\theta(t(b))^{-1}\tau(u^{-1})v_i,
       \qquad b\in I(c).
 \tag{OC6}
\]
Neatness makes this well-defined; the identities for \(t(b)\) give
the required equivariances.  The final identity, together with
\(\theta(s)\equiv1\pmod{\mathfrak m^k}\), proves the asserted
\(N^-_c\)-invariance.  Choosing a deeper congruence for any higher
precision proves that \(s_k(v)\) belongs to \(S(U)\), rather than
merely defining a function at precision \(k\).

The elementary \(U_z\) contraction is
\[
 \begin{split}
 v\ \hbox{fixed by }N,T_c,N^-_{d+1},\quad d\geq c
 \quad\Longrightarrow\quad
 U_zv\ \hbox{fixed by }N,T_c,N^-_d .
 \end{split}
 \tag{OC7}
\]
It follows by writing the double-coset sum and moving an element of
\(N^-_d\) through its representatives; the Iwahori interchange places
the error in \(T_cN^-_{d+1}\).
Apply (OC7) modulo \(\mathfrak m^k\).  For a given \(f\), uniform
continuity supplies an initial opposite-unipotent depth, and (OC1)
makes \(T_c\) act trivially.  Thus, for all sufficiently large \(a\),
\(T_z^af\) is \(N^-_c\)-invariant modulo \(\mathfrak m^k\).
The difference
\[
 T_z^af-s_kE_k(T_z^af)
 \tag{OC8}
\]
has zero values at the representatives and is invariant modulo
\(\mathfrak m^k\) under the full level.  Hence (OC8) belongs to \(V_k\).
Also \(E_k(V_j)\subset(\mathfrak m^j)^{bm}\).

The finite set \((A/\mathfrak m^k)^{bm}\) implies that two sufficiently
late evaluations \(E_k(T_z^af)\) coincide modulo \(\mathfrak m^k\).
Equation (OC8) then makes the \(T_z\)-orbit eventually periodic
modulo \(V_k\).  Since \(T_zV_k\subset V_k\), all subsequent iterates
have the same period.  Consequently \(T_z^{r!}f\) is Cauchy modulo
every \(V_k\).
The valuewise limit belongs to \(S(U)\): left and right equivariance
pass to limits, and at each fixed precision a sufficiently late
approximating function supplies the required continuity depth.
Separatedness and linearity give the map \(e_U\) in (OC4).
Its image is exactly (OC2), it fixes that image, and it preserves
every \(V_k\).

We next prove finite freeness, without assuming it while constructing
the projector.  Applying \(e_U\) to (OC8), and using ordinary
recurrence, gives for every \(k\) finite maps \(E_k,s_k\) with
\[
 f-e_Us_kE_kf\in V_k\qquad(f\in O(U)).
 \tag{OC9}
\]
For precision \(1\), write
\(\sigma=e_Us_1:A^b\to O(U)\) and
\(E=E_1|_{O(U)}\), enlarging \(b\) to include coefficient coordinates.
If \(f\in O(U)\cap V_j\), use (OC9) at precision \(j+1\).
Its coefficients lie in \(\mathfrak m^j\).
Apply (OC9) at precision \(1\) to each of the finitely many ordinary
vectors \(e_Us_{j+1}(e_i)\), multiply the errors by those coefficients,
and sum.  This gives
\[
 \exists a\in(\mathfrak m^j)^b,\qquad
 f-\sigma a\in V_{j+1}.
 \tag{OC10}
\]
Iteration starting at \(j=0\), followed by completeness of the finite
module \(A^b\), proves that \(\sigma\) is surjective.
Starting at a general \(j\) proves
\(O(U)\cap V_j\subset\mathfrak m^jO(U)\); the reverse inclusion follows
directly from values.  This proves (OC5).

Now \(C=\sigma E\) is an endomorphism of the finite \(A\)-module
\(O(U)\) congruent to the identity modulo \(\mathfrak mO(U)\).
Nakayama makes \(C\) surjective, and a surjective endomorphism of a
finite module over a Noetherian ring is bijective.
Thus \(EC^{-1}\) is a section of \(\sigma\), so \(O(U)\) is finite
projective over the local ring \(A\), hence finite free.

Finally, ordinary recurrence modulo \(V_1=\mathfrak mO(U)\) says that
\(T_z\) is surjective modulo \(\mathfrak m\), so the same argument makes
\(T_z\) invertible on \(O(U)\).
For any positive diagonal \(a\), choose \(r\) so that \(z^ra^{-1}\)
is positive.  The positive \(U\)-operators commute and multiply
according to the positive monoid.  Since their product
\(U_{z^ra^{-1}}U_a=U_z^r\) is invertible, \(U_a\) is invertible.
\end{proof}

\begin{srlemma}[Auxiliary group-ring projectivity and exact norm control]
\label{lem:ordinary-aux-projective}
Let \(V\triangleleft U\) have finite index and set \(\Delta=U/V\).
Then \(O(V)\) is finite projective over \(A[\Delta]\).  Moreover the
unnormalized norm
\[
 \operatorname{Nm}_{U/V}=\sum_{\delta\in\Delta}\delta
\]
has image \(O(U)\), and
\[
 \ker\bigl(\operatorname{Nm}_{U/V}:O(V)\to O(U)\bigr)
       =\sum_{\delta\in\Delta}(\delta-1)O(V).
 \tag{OC11}
\]
In particular \(O(V)_\Delta\simeq O(U)\).
No invertibility of \(|\Delta|\) in \(A\) is required.
\end{srlemma}

\begin{proof}
Choose one representative \(R(g)\) for each double coset at level
\(U\times I(0)\).  In
\(g=\gamma R(g)(u,b)\), neatness makes the \(U\)-coordinate \(u\)
unique.  Let \(\chi(g)\) be \(1\) when \(u\in V\), and \(0\) otherwise.
Then \(\chi\) is left \(\Gamma\)-invariant and right
\(V\times I(0)\)-invariant, and
\[
 \sum_{\delta\in\Delta}\chi(g\delta)=1.
 \tag{OC12}
\]
Indeed precisely one right coset sends the \(U\)-coordinate to \(V\).
Multiplication by \(\chi\) maps \(O(V)\) into \(S(V)\), since \(\chi\)
is \(I(0)\)-invariant.  Put
\(\Phi=e_V\circ M_\chi\in\operatorname{End}_A(O(V))\).
The ordinary projector commutes with the away translations: this
holds for each \(U_z^{r!}\), hence for their separated limits.
Taking \(e_V\) of the partition identity (OC12) yields
\[
 \sum_{\delta\in\Delta}\delta\Phi\delta^{-1}=1.
 \tag{OC13}
\]
Since \(O(V)\) is finite free over \(A\), the equivariant surjection
\[
 A[\Delta]\otimes_AO(V)\longrightarrow O(V),\qquad
 \delta\otimes v\longmapsto\delta v
\]
has the \(A[\Delta]\)-linear section
\[
 v\longmapsto
 \sum_{\delta\in\Delta}\delta\otimes\Phi(\delta^{-1}v).
\]
Linearity follows by reindexing and the composite is the identity by
(OC13).  This proves finite projectivity.

The fixed submodule \(O(V)^\Delta\) is exactly \(O(U)\), by the
function definitions and the level-independent recurrence (OC2).
For \(v\) fixed, (OC13) gives
\(\operatorname{Nm}(\Phi(v))=v\), proving norm surjectivity.
If \(\operatorname{Nm}(v)=0\), then
\[
 \sum_\delta\Phi(\delta^{-1}v)=\Phi(\operatorname{Nm}(v))=0
\]
and (OC13) gives
\[
 v=\sum_\delta
       \bigl(\delta\Phi(\delta^{-1}v)-\Phi(\delta^{-1}v)\bigr).
\]
This proves the nontrivial inclusion in (OC11); the reverse inclusion
is the norm identity.
\end{proof}

\paragraph{Hecke localization and the actual smooth auxiliary module.}
Use only the tame Hecke operators away from the varying auxiliary
places, together with the ordinary operators at \(\ell\).
They commute on \(O(U)\); their image
\(\mathbb T_U\subset\operatorname{End}_A(O(U))\) is a finite commutative
\(A\)-algebra.  Completeness of \(A\) decomposes \(\mathbb T_U\) into
its local factors.  The factor selected by the prescribed residual
Hecke character is a central-idempotent summand
\[
 M(U)=\epsilon_UO(U).
 \tag{OC14}
\]
It is finite free over \(A\).  At level \(V\triangleleft U\), the
projector \(\epsilon_V\) commutes with the \(\Delta\)-action, because
every selected Hecke operator does.  Thus \(M(V)\) is a projective
\(A[\Delta]\)-summand of \(O(V)\).

For clarity about invariants, view all \(M(K)\), for compact open
\(K\) in the product of the varying local groups, in the single
function space, and put
\[
 \mathcal M=\bigcup_KM(K).
 \tag{OC15}
\]
This union is directed and translation-stable, so it is an actual
smooth auxiliary representation.  It satisfies
\(\mathcal M^K=M(K)\).
Here the only point beyond the function invariance is compatibility
of the localized summand with intersection.  If \(Y\subset Y'\) are
the corresponding finite ordinary modules, then Artin--Rees gives a
constant \(c\) with
\(Y\cap\mathfrak m^{r+c}Y'\subset\mathfrak m^rY\).
Consequently the simultaneous generalized residual-eigenvalue
condition, expressed at every \(\mathfrak m\)-adic precision, cuts out
the same submodule of \(Y\) whether tested in \(Y\) or \(Y'\).
Intersecting two levels and using this observation proves the claimed
fixed-vector equality.
The quotient \(K/K'\) acting on \(M(K')\) is the actual auxiliary
quotient: the local embedding used to form the tame level has
preimage \(K'\) inside \(K\).  Lemma
\ref{lem:ordinary-aux-projective} and the commuting central idempotent
therefore prove that \(\mathcal M^{K'}\) is projective over
\(A[K/K']\).

These constructions supply finite freeness, group-ring projectivity,
the smooth module, and the norm/augmentation control in the
finite-level input.  They do not by themselves prove that the chosen
localized factor is nonzero, that its Hecke determinant is realized
by a Galois representation, or that every classical specialization
has the prescribed local-global compatibility.  Those are separate
automorphic arithmetic inputs.  In particular this construction must
not be used to declare the automorphic seed proved.

\paragraph{Exact source locations.}
The preceding construction follows the proofs, rather than just the
output interfaces, in the frozen source
\srcid{alpoge-source-7cbf7a0c}, under
\srcid{Development.AutomorphicArithmetic.AutomorphicVarieties}.
The arithmetic smallness proof is
\srcid{HeckeAlgebras.OrdinaryHeckeCharacters.NtExpand.neatU}.
The interpolation and contraction proofs are
\srcid{BuildingsResidueChars.ExistsLinearSectionAutomorphicFunctionsOfFree},
\srcid{ExistsIwahoriTorusProjection},
\srcid{ExistsSectionAutomorphicSpaceWithPrescribedValue}, and
\srcid{UOperatorPowEventuallyInvariantModPow}.
The algebra proof of (OC4)--(OC5) is
\srcid{HeckeAlgebras.OrdinaryHeckeCharacters.OrdFree.ordAbs},
\srcid{ordStruct}, and \srcid{ordFree}; (OC12)--(OC13) and
group-ring projectivity are proved in
\srcid{OrdProj.ordProj\_chi}, \srcid{ordProj\_phi},
\srcid{ordProj\_higman}, and \srcid{ordProj}.
Hecke localization is \srcid{HiComm.hiComm}; the fixed-level
comparison and localized projectivity are the explicit
\srcid{pb\_locInf} and \srcid{pb\_main} arguments in
\srcid{PiBasic}, culminating in \srcid{piBasic}.
All these file paths are below
\url{https://github.com/ChristianKudeba/alpoge-rh-verification/tree/alpoge-source-7cbf7a0c/src/Development/AutomorphicArithmetic/AutomorphicVarieties}.


\subsection{Base change and classical specialization of the constructed ordinary modules}
\label{sec:ordinary-classical-control}

We continue with the function modules of
\ref{sec:ordinary-module-construction}.  This subsection proves the
ordinary base-change and classical control statements used to obtain
separating specializations.  It does not assume an ordinary
automorphy-lifting theorem.

\begin{srlemma}[Ordinary base change]
\label{lem:ordinary-base-change}
Let \(A\to A'\) be a local homomorphism of complete Noetherian local
rings with finite residue fields.  Base-change the tame coefficient
action and the torus character.  At each of the neat finite levels,
coefficient change induces an isomorphism
\[
 A'\otimes_A O_A(U)\ \xrightarrow{\ \sim\ }\ O_{A'}(U).
 \tag{CC1}
\]
This assertion does not require \(A'\) to be flat over \(A\).
\end{srlemma}

\begin{proof}
Coefficient change carries value filtration \(V_k\) into \(V'_k\).
It intertwines every \(U_z\)-power and hence the factorial ordinary
projectors.  In particular it defines a map \(j:O_A(U)\to O_{A'}(U)\).

Choose one finite interpolation system at precision \(1\), with
evaluation functionals \(E_p\), interpolation sections \(s_p\), and
\(w_p=e_As_p\).  The proof of (OC9), using this fixed interpolation
system, gives
\[
 f-\sum_pE_p(f)w_p\in\mathfrak m_AO_A(U).
 \tag{CC2}
\]
The same sections, changed to \(A'\), give
\[
 f'-\sum_pE'_p(f')j(w_p)\in\mathfrak m_{A'}O_{A'}(U).
 \tag{CC3}
\]
To justify using precisely the same system, both characters are
trivial modulo their maximal ideals on the chosen \(T_c\).
For an ordinary \(f'\), a sufficiently large factorial power is
congruent to \(f'\) modulo \(V'_1\) and is \(N^-_c\)-invariant there,
by (OC7).  Its interpolation has the same representative values.
Subtract, apply that factorial power to the finite interpolation
sections, and use \(e_As_p=w_p\) modulo \(V_1\).
This proves (CC3) with the displayed \(j(w_p)\), rather than with an
independently chosen generating family.

Let \(C(f)=\sum_pE_p(f)w_p\).  By (CC2), Nakayama and the Noetherian
endomorphism lemma make \(C\) an automorphism of \(O_A(U)\).
Let
\[
 L:A'\otimes_AO_A(U)\to O_{A'}(U),\qquad
 R(f')=\sum_pE'_p(f')\otimes w_p.
\]
Evaluation is compatible with \(j\), so
\(RL=1\otimes C\).
The latter is invertible, proving injectivity of \(L\).
Equation (CC3) and Nakayama prove surjectivity of \(L\), since
\(O_{A'}(U)\) is finite over \(A'\).
\end{proof}

\paragraph{The elementary algebraic weight data.}
Let \(\mathcal O\) be a sufficiently large finite coefficient DVR
and let \(\Omega\) be an algebraically closed characteristic-zero
field containing its fraction field.
An algebraic dominant weight at the split \(\ell\)-adic places gives
a finite-dimensional representation \(\xi\) with a lowest-weight
coordinate numbered \(0\), diagonal character \(\chi\), and an
integral weight lattice.  Enlarge \(\mathcal O\) to contain the
finitely many coefficient-field embeddings used to define it.
The properties needed below are the following concrete properties
of this weight lattice:
\[
 \begin{array}{ll}
 \text{(i)}&
  e_0^\vee\xi(n)=e_0^\vee\quad(n\in N),\qquad
  e_0^\vee\xi(t)=\chi(t)e_0^\vee\quad(t\in T);\\[2mm]
 \text{(ii)}&
  \xi(N)\ \hbox{is integral},\qquad
  \chi(a)^{-1}\xi(a)\ \hbox{is integral for }a\hbox{ positive};\\[2mm]
 \text{(iii)}&
  \chi(z)^{-1}\xi(z)\ \hbox{has every column except column }0
                     \hbox{ in }\mathfrak m_{\mathcal O};\\[2mm]
 \text{(iv)}&
  \xi(y)\equiv1\pmod{\mathfrak m_{\mathcal O}^k}
   \quad(y\in N^-_c,\ c\hbox{ sufficiently large for }k).
 \end{array}
 \tag{CC4}
\]
One may construct the lattice using the integral Schur modules for
the successive differences of the highest weight, tensoring with
a determinant power.  The matrix coefficients of unipotents are
integral polynomials, which proves (i), (ii), and (iv).
In a weight basis the normalized action of the strictly positive
\(z\) is \(1\) at the extremal coordinate and has positive valuation
at every other weight; its extremal weight has multiplicity one.
This proves (iii).
These observations work for arbitrary \(\ell\); no
Fontaine--Laffaille range is used in this representation-theoretic
control argument.

Suppose now that the specialized weight character has the form
\[
 \theta(t)=\phi(t)\chi(t),\qquad \phi|_{T_c}=1
 \tag{CC5}
\]
for a finite-depth character \(\phi:T\to\mathcal O^\times\).
Choose \(c\) sufficiently deep for its finite part and the integral
weight representation.
Let \(C(U,c,\phi)\) be the classical algebraic function space:
its functions \(F\) transform on the left by \(\xi\) at the
\(\ell\)-adic component of \(\Gamma\), satisfy the tame coefficient
transformation, and have right Iwahori character \(\phi^{-1}\)
at depth \(c\).
For these functions set
\[
 qF(g)=\xi(g_\ell)^{-1}F(g),\qquad
 \operatorname{ctrl}(F)(g)=(qF(g))_0 .
 \tag{CC6}
\]
The coefficient index from the tame lattice is carried along in
this formula.  The classical ordinary operator is
\[
 U_z^{\mathrm{cl}}=
   \chi(z)^{-1}\sum_{n\in N/(N\cap zNz^{-1})}[nz].
 \tag{CC7}
\]
The row identities in (CC4)(i) give, by evaluating the finite sum,
\[
 \operatorname{ctrl}(U_a^{\mathrm{cl}}F)
        =U_a\,\operatorname{ctrl}(F)
 \tag{CC8}
\]
for every positive \(a\).  The control map also intertwines away
translations and finite away-Hecke sums.

Evaluation of all coordinates of \(qF\) at representatives for
\(\Gamma\backslash(G_{\mathrm a}\times G_\ell)/(U\times I(c))\)
identifies \(C(U,c,\phi)\) with a finite-dimensional coordinate
space over \(\Omega\).
Indeed equivariance recovers the function from these coordinates;
arbitrary coordinates give a function because the level is neat.
The integral coordinate lattice \(L_c\) consists exactly of those
functions for which every coordinate of \(qF(g)\) belongs to
\(\mathcal O\) for every \(g\).
The integral Iwahori action shows the equivalence with integrality
at the chosen representatives.
Consequently \(L_c\) is finite free and is preserved by (CC7).

\begin{srproposition}[Classical ordinary control]
\label{prop:classical-ordinary-control}
Let \(C(U,c,\phi)^{\mathrm{ord}}\) be the sum of generalized
eigenspaces of \(U_z^{\mathrm{cl}}\) whose eigenvalue is an integral
unit.  Then \(\operatorname{ctrl}\) is injective on this space and
\[
 \operatorname{ctrl}\bigl(C(U,c,\phi)^{\mathrm{ord}}\bigr)
  =\Omega\cdot O_{\mathcal O}(U)
 \tag{CC9}
\]
inside the common scalar-valued function space.
If \(A\to\mathcal O\to\Omega\) is an arithmetic specialization,
(CC1) and (CC9) therefore give a specialization map from the
constructed \(A\)-ordinary module to classical ordinary forms.
The images of any \(A\)-basis are linearly independent over
\(\Omega\).
The map commutes with all selected Hecke operators and all away
translations, including the auxiliary local group actions.
\end{srproposition}

\begin{proof}
The ordinary part of the finite integral lattice \(L_c\) is the
image of \(\lim_r(U_z^{\mathrm{cl}})^{r!}\).
This follows from the same finite-quotient periodicity proof as
(OC4), now on a finite free module.
After extension to \(\Omega\), its span is exactly the direct sum
of the generalized eigenspaces with integral-unit eigenvalue:
the complementary eigenvalues have positive valuation, and the
factorial projector is \(0\) there and \(1\) on the unit part.

The contraction controlling the kernel is particularly simple.
If every coordinate of \(qF\) belongs to \(\mathfrak m^j\), while
its coordinate \(0\) belongs to \(\mathfrak m^{j+1}\), then
\[
 q(U_z^{\mathrm{cl}}F)\ \hbox{has all coordinates in }
                   \mathfrak m^{j+1}.
 \tag{CC10}
\]
For an individual summand of (CC7), its untwisted value is
\[
 \xi(n)\bigl(\chi(z)^{-1}\xi(z)\bigr)
                 qF(g(nz)).
\]
Columns different from \(0\) contribute
\(\mathfrak m\mathfrak m^j\), by (CC4)(iii).
Column \(0\) contributes \(\mathfrak m^{j+1}\) by hypothesis.
The integral matrix \(\xi(n)\) and the finite sum preserve the
ideal, proving (CC10).

If \(F\) belongs to the integral ordinary lattice and
\(\operatorname{ctrl}(F)=0\), (CC10), ordinary recurrence modulo
every \(\mathfrak m^kL_c\), and stability of these lattices show
inductively that \(F\in\mathfrak m^jL_c\) for every \(j\).
Separatedness gives \(F=0\).
The same proves injectivity after scalar extension: for finitely
many integral independent functions, independence is preserved
when the fraction field is extended to \(\Omega\).
To verify this last elementary fact, clear denominators first over
\(\operatorname{Frac}\mathcal O\), then expand the finitely many
extension-field coefficients in a basis of their span over that
fraction field and compare coefficients in the function identity.
It follows that the control map is injective on the whole classical
ordinary space.

For an integral ordinary classical \(F\), the function
\(f=\operatorname{ctrl}(F)\) has integral values.
It satisfies the scalar equivariances with \(\theta\) by
(CC4)(i) and (CC5).
It is uniformly continuous along the opposite unipotent:
\(\xi(y)\equiv1\pmod{\mathfrak m^k}\) at sufficiently deep level,
by (CC4)(iv).
Equation (CC8) transfers ordinary recurrence to \(f\).
Thus the control image is contained in \(\Omega\cdot O_{\mathcal O}(U)\).

We prove the reverse inclusion by an explicit approximation.
Take \(f\in O_{\mathcal O}(U)\).
Choose \(c'\geq c\) so deep that both \(f\) and \(\xi\) are
\(N^-_{c'}\)-invariant modulo \(\mathfrak m\).
At the finite representatives at depth \(c'\), prescribe the
coordinates
\[
 qF(x)_0=f(x),\qquad qF(x)_i=0\quad(i\ne0).
\]
Neat interpolation gives an integral classical form at that depth.
Its scalar control differs from \(f\) by valuewise
\(\mathfrak m\), since the two controls have the same
representative values and the same congruence invariance.

A sufficiently large positive power of \(U_z^{\mathrm{cl}}\)
moves this form from depth \(c'\) to depth \(c\).
This is the same one-step Iwahori interchange as (OC7):
the opposite-unipotent depth drops by one at each application;
the finite character is already trivial on \(T_c\).
Choose the power also divisible by a recurrence period for \(f\)
modulo \(\mathfrak m\).
By (CC8), its scalar control still agrees with \(f\) modulo
\(\mathfrak m\).
Apply the classical factorial ordinary projector at depth \(c\).
The resulting ordinary classical form has control congruent to
\(f\), because (CC8) identifies its control limit with the scalar
ordinary projector, which preserves \(V_1\) and fixes \(f\).

Let
\[
 L=\bigl(\operatorname{ctrl}
       (C(U,c,\phi)^{\mathrm{ord}})\bigr)
               \cap O_{\mathcal O}(U).
\]
This is an \(\mathcal O\)-submodule of a finite
\(\mathcal O\)-module.  The preceding approximation proves
\(L+\mathfrak m O_{\mathcal O}(U)=O_{\mathcal O}(U)\);
here valuewise \(\mathfrak m\) is
\(\mathfrak mO_{\mathcal O}(U)\) by (OC5).
Nakayama gives \(L=O_{\mathcal O}(U)\), proving (CC9).

Finally apply (CC1).
A basis of \(O_A(U)\) becomes an \(\mathcal O\)-basis of
\(O_{\mathcal O}(U)\).
As integral scalar-valued functions these basis vectors remain
independent over \(\Omega\), by the coefficient comparison just
used.  Their unique inverse-control images are therefore
independent classical forms.
All constructions commute with the indicated Hecke operators:
coefficient change and control do so by their pointwise formulas,
and the ordinary projectors do so by their factorial limits.
The same formulas commute with away translations, which proves
the asserted auxiliary equivariance.
\end{proof}

\paragraph{Automatic continuity of arithmetic coefficient maps.}
The use of a local specialization in (CC1) is justified
intrinsically.  Let \(T\) be a complete Noetherian local
\(\mathcal O\)-algebra with the same finite residue field and let
\(x:T\to\overline{\mathbb Q}_\ell\) be an
\(\mathcal O\)-algebra map.
For \(a\in\mathfrak m_T\) one has \(|x(a)|_\ell<1\).
Otherwise \(x(a)^{-1}\) is integral over \(\mathbb Z_\ell\);
reversing its monic integral equation gives a polynomial
\(1+c_1X+\cdots+c_rX^r\) vanishing at \(x(a)\), with
\(c_i\in\mathbb Z_\ell\).
Its value at \(a\) is a unit of \(T\), a contradiction.
Writing an arbitrary element as its \(\mathcal O\)-residue
representative plus an element of \(\mathfrak m_T\) gives
\(|x(t)|_\ell\leq1\) for every \(t\).
For finitely many generators of \(\mathfrak m_T\), take the maximum
of their norms to obtain \(b<1\); then
\[
 |x(\mathfrak m_T^j)|_\ell\leq b^j.
\]
Thus \(x\) is continuous.  Choose finitely many topological
generators \(t_1,\ldots,t_s\) of \(T\) over \(\mathcal O\); these
exist by completeness, finite generation of \(\mathfrak m_T\), and
surjectivity on the residue field.  The finitely many algebraic
elements \(x(t_i)\), together with the coefficient field, belong to
one finite extension \(K/\mathbb Q_\ell\).
Every element of \(T\) is an adic limit of polynomials in the
\(t_i\); continuity and completeness of \(K\) imply
\(x(T)\subset\mathcal O_K\).
The map \(T\to\mathcal O_K\) is local, since maximal-ideal elements
have norm less than one.
This gives the finite coefficient DVR required in (CC1).
The residue-preserving complete local coefficient rings in the
application therefore have exactly the continuity and coefficient
field containment needed for ordinary specialization; continuity is
not silently imposed on an arbitrary field-valued algebra map.

\paragraph{Remaining arithmetic input and exact sources.}
This proves the module-theoretic classical specialization and
basis independence.  To pass from classical forms to attached
Galois representations, one still uses the definite-unitary
automorphic decomposition, its attached semisimple Galois
representations, and local-global compatibility.  To separate
Hecke-algebra elements by these specializations one additionally
uses density of the algebraic weight characters in the specified
weight algebra.  Neither the Galois-realization assertion nor
weight density is being inferred solely from the control
isomorphism.

The exact frozen source is under
\srcid{Development.AutomorphicArithmetic.AutomorphicVarieties},
tag \srcid{alpoge-source-7cbf7a0c}.
Base change is the finite evaluation argument
\srcid{HeckeAlgebras.OrdinaryHeckeCharacters.OrdBCMapOrd.ordBC\_abs},
\srcid{ordBC\_key}, and \srcid{ordBC}.
Control is
\srcid{CtrlOrd.ctrlOrd\_inj}, \srcid{ctrlOrd\_approx},
and \srcid{ctrlOrd}; the elementary lattice calculation (CC10) is
\srcid{BuildingsResidueChars.UpOperatorMapsLatticeOfControlValues}.
The other explicit construction helpers are
\srcid{ExistsIntegralFormLiftingControlMap},
\srcid{ExistsAutomorphicFormWithPrescribedIwahoriValues},
and \srcid{ExistsFreeLatticeSpanningOrdinaryPart}.
The coefficient-change independence lemma is
\srcid{HeckeAlgebras.OrdinaryHeckeCharacters.FnIndep.fnIndep}.
The simultaneous specialization and all Hecke/auxiliary
equivariances are assembled by the actual proof
\srcid{HeckeAlgebras.HeckeFixedLevelSpaces.FreeMap.absSpec},
then by \srcid{OrdinaryHeckeCharacters.SpecIntIff.spec}.
The continuity calculation is the proof
\srcid{DenseThresholdCases.TCoreDet.cd\_norm\_lt},
\srcid{cd\_pow}, and \srcid{cd\_cont}.
All files are below
\url{https://github.com/ChristianKudeba/alpoge-rh-verification/tree/alpoge-source-7cbf7a0c/src/Development/AutomorphicArithmetic/AutomorphicVarieties}.


\subsection{Explicit old/new maps from the integral Iwahori algebra}
\label{sec:explicit-old-new}

We construct the two maps required in the auxiliary control argument.
They are defined by finite sums of Hecke operators.  In particular the
construction never divides by the index of an Iwahori subgroup in a
maximal compact subgroup.  The only arithmetic coefficient comparison
used in this subsection is displayed in (ON8) below.

\paragraph{The integral calculation at one split place.}
Let \(k\) be a nonarchimedean local field, with residue cardinality
\(Q\), \(K=\mathrm{GL}_n(\mathcal O_k)\), and standard Iwahori \(I\).
Let \(Q\) be a unit in the coefficient ring.
Write \(X_i\), for \(0\leq i<n\), for the commuting Bernstein operators
with the normalization
\[
 X_i=Q^i\,\Theta(\operatorname{diag}_i(\varpi)).
 \tag{ON1}
\]
Thus, if \(t_r=\operatorname{diag}(\varpi,\ldots,\varpi,1,\ldots,1)\)
has \(r\) entries \(\varpi\), then
\[
 X_i[I t_i I]=Q^i[I t_{i+1}I].
\]
The positive diagonal Hecke operators are units when \(Q\) is a unit.
One can see this in the affine Iwahori algebra: the simple operators
satisfy
\(T_s^2=(Q-1)T_s+Q\), hence
\(T_s^{-1}=Q^{-1}(T_s-(Q-1))\), and translations are products of these
operators and length-zero units.  The commuting positive diagonal
action consequently extends to the torus action \(\Theta\) in (ON1).
This proves existence of the torus action used here directly from
double-coset operators.

Put, in the polynomial ring \(\mathcal P=\mathbb Z[Q,Q^{-1}]
[X_0,\ldots,X_{n-1}]\),
\[
 D(X)=\prod_{i<j}(X_j-X_i),\qquad
 N_Q(X)=\prod_{i<j}(X_j-QX_i),\qquad
 \partial(f)=
 \frac{\sum_{w\in S_n}\operatorname{sgn}(w)\,w(f)}{D(X)}.
 \tag{ON2}
\]
The numerator is divisible by \(D\) over this integral polynomial
ring: it vanishes after each substitution \(X_i=X_j\), so the distinct
prime linear factors divide it.  In particular no factorial is
inverted in \(\partial\).  Its values are symmetric polynomials and
\(\partial(sf)=s\partial(f)\) for symmetric \(s\).

Let \(E=[K:I]\) denote the trace operator from Iwahori-fixed to
\(K\)-fixed vectors, considered as an endomorphism of the
Iwahori-fixed module.  This notation denotes the double-coset sum,
not scalar multiplication by the integer \([K:I]\) on arbitrary
vectors.  In the Iwahori algebra \(E=\sum_{w\in S_n}T_w\).
Its integral identities are
\[
 \begin{split}
 Es(X)&=s(X)E &&(s\ \hbox{symmetric}),\\
 E f(X)E&=\partial(N_Qf)(X)\,E .
 \end{split}
 \tag{ON3}
\]
Here and below a polynomial on the right means its Bernstein action.

For completeness, these identities can be verified in the universal
Bernstein algebra.  Its generators satisfy the finite Hecke
quadratic and braid relations, and, for adjacent \(i,j=i+1\),
\[
 \begin{split}
 T_iX_i&=X_jT_i+(Q-1)X_i,\\
 T_iX_j&=X_iT_i-(Q-1)X_i,\qquad
 T_iX_a=X_aT_i\quad(a\ne i,j).
 \end{split}
 \tag{ON4}
\]
These are the rank-one double-coset calculation with
\(I\operatorname{diag}(\varpi,1)I\); operators at other coordinates
commute.  To check the universal identities without making a
specialization injective, use the fraction field of \(\mathcal P\)
and the crossed product by \(S_n\).  There
\[
 T_i=
 \frac{(1-Q)X_i}{X_j-X_i}
   +\frac{QX_j-X_i}{X_j-X_i}\,[s_i].
 \tag{ON5}
\]
These operators satisfy (ON4), the quadratic relation, and the braid
relation, by the two-variable and three-variable expansions,
respectively.  Straightening (ON4) spans the universal algebra by
\(\mathcal P T_w\).  In (ON5), the coefficient of \([w]\) in \(T_w\)
is a nonzero product over the inversions of \(w\); all remaining
terms have smaller Bruhat length.  Thus these spanning elements are
independent and the crossed-product calculation is faithful.

The equality \(T_iE=QE\) implies that if \(A_w\) is the coefficient of
\([w]\) in \(E\), then \(A_{s_iw}=s_i(A_w)\).
The coefficient at the longest permutation is obtained solely from
its longest reduced word, and is
\(w_0(N_Q/D)\).  Hence
\[
 E=\sum_w w(N_Q/D)[w],\qquad
 E(f)=\partial(N_Qf).
 \tag{ON6}
\]
It follows that \(E\) commutes with symmetric polynomials and that
\(E f E=\partial(N_Qf)E\), since \(E\) has symmetric image.
The result is an identity in the integral universal algebra, by
faithfulness; (ON5) is only a verification device.
Specializing \(Q\) and applying the actual Iwahori action proves
(ON3) over the original coefficient ring.

Here is an explicit finite divided-difference dual family.
Let \(\mathcal R=\{(i,j):0\leq i<j<n\}\).
For each subset \(H\subset\mathcal R\), put
\[
 a_H=\prod_{(i,j)\in H}X_j,\qquad
 c_H=\prod_{(i,j)\in\mathcal R\setminus H}(-X_i).
\]
Then
\[
 \sum_H a_H\,w(c_H)=
 \prod_{i<j}(X_j-X_{w(i)})
 =\begin{cases}D,&w=1,\\0,&w\ne1.\end{cases}
\]
For \(w\ne1\), choose \(i\) with \(i<w(i)\); the factor indexed by
\((i,w(i))\) is zero.  Expanding the antisymmetrizer therefore gives
\[
 \sum_Ha_H\partial(c_Hf)=f,\qquad
 \sum_Hc_H\partial(a_Hf)=f,\qquad
 \sum_H a_H(X)E c_H(X)=N_Q(X).
 \tag{ON7}
\]
The second equality follows by applying a permutation to the dual
identity for its inverse, which interchanges the two families.
For the last equality, insert (ON6); only the identity permutation
survives in its crossed-product expansion.
These formulas prove the precise integral identity used in the
old/new construction.

\paragraph{The global modules and the coefficient comparison.}
Let \(\mathcal Q=\{w_1,\ldots,w_q\}\) be the distinct split auxiliary
places, away from \(\ell\) and from the fixed ramification set.
Use the smooth module \(\mathcal M\) constructed in
\ref{sec:ordinary-module-construction}; at these places write
\[
 V=\mathcal M^{\prod_j I_j},\qquad
 M=\mathcal M^{\prod_j K_j}.
\]
The spherical module is identified with the fixed base module by the
actual maximal-level comparison.  More generally its coefficients
are related by the ring map \(\beta:R\to B\); the inclusion
\(\iota:M\to V\) is additive and satisfies
\(\iota(\beta(r)m)=r\iota(m)\).
Let
\[
 P_j(T)\in B[T]
\]
be monic of degree \(n\).
The comparison with the base Galois determinant is exactly the
following spherical coefficient identity:
\[
 e_k(U_{j,0},\ldots,U_{j,n-1})\iota(m)
   =\iota\bigl((-1)^k[P_j]_{n-k}\,m\bigr),
 \quad 0\leq k\leq n.
 \tag{ON8}
\]
The normalization of \(U_{j,i}\) is (ON1).
In the application
\[
 P_j(T)=D_{\mathrm{base}}(T-[\mathrm{Frob}_{w_j}]),
\]
with the inverse-Frobenius convention of the original
\srcid{Pair}; (ON8) uses that same convention.

We spell out the proof of (ON8) from the arithmetic realization
input, so that it is not an additional lifting assumption.
The actual classical specializations separate vectors of
\(\mathcal M\).  Under such a specialization the auxiliary action
is an isotypic sum of an irreducible smooth complex local
representation \(\sigma_j\).  A nonzero image of a spherical vector
forces \(\sigma_j^{K_j}\ne0\).
The normalized Bernstein eigenvalues \(\alpha_i\) of this spherical
representation satisfy
\[
 \operatorname{charpoly}(\rho(\mathrm{Frob}_{w_j}))
       =\prod_i(T-\alpha_i)
\]
by the unramified local-global comparison with the attached Galois
representation.  The symmetric Bernstein operator \(e_k\) acts on
the spherical line by \(e_k(\alpha)\); Vieta's formula identifies
this with the specialized right hand side of (ON8).
If the specialized spherical vector is zero both sides are zero.
Separation by all these specializations proves (ON8) integrally.
The argument needs precisely the separating arithmetic
specializations, attached determinant, and this unramified
coefficient comparison; it does not need a theorem lifting
individually polarized residual constituents.

Let \(E_j:V\to V\) be the unnormalized local Iwahori-to-spherical
trace, and put \(E=\prod_jE_j\).
The operators at distinct places commute and \(E(V)\subset M\).
Write \(\operatorname{tr}:V\to M\) for this product trace under the
base identification, so \(\iota\operatorname{tr}=E\).
The module \(\mathcal M\) has no nonzero integer torsion: for example
this follows from its separating complex specializations, or from
its finite free weight-algebra fixed levels.
For \(v\in V\) spherical at \(w_j\), the sandwich identity in (ON3)
gives
\[
 E_j f(U_j)v=\partial(N_{Q_j}f)(U_j)v.
 \tag{ON9}
\]
Indeed \(E_jv=[K_j:I_j]v\); first multiply the proposed identity by
this positive integer and use (ON3), then cancel the integer using
torsion freeness.  Thus no assertion that \([K_j:I_j]\) is a unit
is hidden in (ON9).
The product trace is \(\beta\)-semilinear, and, by (ON8) and
symmetric commutation, it intertwines every symmetric polynomial
with its image in the coefficients of \(P_j\).

\begin{srproposition}[The two concrete maps]
\label{prop:concrete-two-maps}
Let
\[
 S=B[u_{j,i}:j,i]\big/
 \left(\hbox{coefficients of }
          \prod_i(T-u_{j,i})-P_j(T)\ \hbox{for every }j\right)
\]
be the simultaneous ordered splitting algebra, and let
\(\overline f\) denote coefficient change by \(\beta\) followed by
the quotient map to \(S\).
There are additive maps, intertwining the polynomial actions,
\[
 \Phi:S\otimes_BM\longrightarrow V,\qquad
 \Psi:V\longrightarrow S\otimes_BM,
\]
such that, for
\[
 N_{\mathcal Q}(u)=
     \prod_{j=1}^q\prod_{i<i'}(u_{j,i'}-Q_ju_{j,i}),
\]
\[
 \Phi\Psi=N_{\mathcal Q}(U),\qquad
 \Psi\Phi=\overline{N_{\mathcal Q}}.
 \tag{ON10}
\]
They may be chosen so that the two composites are instead
multiplication by
\[
 f_{\mathcal Q}(u)=
 \prod_j\prod_i\prod_{i'\ne i}
       (u_{j,i}-u_{j,i'})(u_{j,i}-Q_ju_{j,i'}).
 \tag{ON11}
\]
These are the two maps appearing in the finite-level ordinary
control relation.
\end{srproposition}

\begin{proof}
Define \(\Phi\) first on polynomials by
\[
 \Phi\!\left(\overline{\sum_\nu b_\nu u^\nu}\otimes m\right)
       =\sum_\nu U^\nu\iota(b_\nu m).
 \tag{ON12}
\]
The coefficient relations defining \(S\) kill this expression by
(ON8) and Vieta's formula.  Thus (ON12) descends to the quotient
and the tensor product.  It satisfies
\(\Phi(1\otimes m)=\iota(m)\) and
\(\Phi(\overline h\,z)=h(U)\Phi(z)\).
This definition is meaningful even when \(\beta\) is not surjective.

For a tuple \(\mathbf H=(H_1,\ldots,H_q)\) of subsets of
\(\mathcal R\), put
\[
 A_{\mathbf H}=\prod_j a_{H_j}(u_j),\qquad
 C_{\mathbf H}=\prod_j c_{H_j}(u_j).
\]
The reverse map is the finite sum
\[
 \Psi(v)=
 \sum_{\mathbf H}
    \overline{A_{\mathbf H}}\otimes
       \operatorname{tr}\bigl(C_{\mathbf H}(U)v\bigr).
 \tag{ON13}
\]
We verify its polynomial equivariance, since this is essential for
the second composite.  In
\(\mathcal P\otimes_{\mathcal P^{S_n}}\mathcal P\), (ON7) gives
\[
 \sum_Ha_H\otimes c_Hx=\sum_Hxa_H\otimes c_H .
 \tag{ON14}
\]
Indeed expand \(c_Hx\) by the second formula of (ON7), transfer
the symmetric coefficient
\(\partial(a_Kc_Hx)\) across the tensor sign, and use the first
formula of (ON7) to sum over \(H\).
After taking the product of (ON14) over the auxiliary places,
symmetric balance is precisely the symmetric intertwining property
of \(\operatorname{tr}\) and (ON8).
Consequently
\(\Psi(h(U)v)=\overline h\,\Psi(v)\).

Applying (ON12) to (ON13) and using
\(\iota\operatorname{tr}=\prod_jE_j\) gives
\[
 \Phi\Psi
 =\sum_{\mathbf H}A_{\mathbf H}(U)
            \left(\prod_jE_j\right)C_{\mathbf H}(U)
 =\prod_jN_{Q_j}(U_j),
\]
where (ON7) and commutation between distinct places give the final
equality.
For a spherical vector, repeated use of (ON9) yields
\[
 \operatorname{tr}\bigl(C_{\mathbf H}(U)\iota(m)\bigr)
  =\left(\prod_j
        \partial(N_{Q_j}c_{H_j})\right)(u)\,m.
\]
The factor on the right is symmetric at every place and hence is a
well-defined \(B\)-coefficient under (ON8).
The first equality of (ON7), applied to \(f=N_{Q_j}\), now gives
\[
 \Psi(\iota(m))=\overline{N_{\mathcal Q}}(1\otimes m).
\]
Since \(S\otimes_BM\) is generated from these elements by the
polynomial action, polynomial equivariance proves the second
identity in (ON10).

Finally \(N_{\mathcal Q}\) divides \(f_{\mathcal Q}\) over the integral
polynomial ring: each factor \(u_{j,i'}-Q_ju_{j,i}\), \(i<i'\),
occurs among the ordered factors in (ON11).
Write \(f_{\mathcal Q}=N_{\mathcal Q}c\).
Keep \(\Phi\) and replace \(\Psi\) by \(\Psi\circ c(U)\).
Equivariance and (ON10) make both composites multiplication by
\(f_{\mathcal Q}\), as claimed.
\end{proof}

\paragraph{What this proves and what still comes from arithmetic.}
The maps (ON12)--(ON13) are now constructed, with exactly the
discriminant used by the auxiliary support argument.  Together with
the finite-projective modules and norm control constructed in the
preceding subsection, and the uniform auxiliary nilpotence proof,
this supplies the module-theoretic finite-level control mechanism.
It does not supply an automorphic seed or construct attached Galois
representations from nothing.
The only remaining arithmetic realization premise in this particular
old/new step is the explicitly described classical separation and
coefficient comparison producing (ON8).
In particular the presence of a two-map interface in the formal
source is not being substituted for its proof.

\paragraph{Exact source locations.}
Use the frozen tag \srcid{alpoge-source-7cbf7a0c}, under
\srcid{Development.AutomorphicArithmetic.AutomorphicVarieties}.
Definitions of the splitting algebra and of the exact two-map
conclusion are
\srcid{AlgebraTools.MSplitDefs} and
\srcid{PatchingModules.NilpotentLevelLowering.LevelModule.TwoMaps}.
The actual normalized local Bernstein operators, their faithfulness,
the trace identities, and the divided-difference formulas are
\srcid{AlgebraTools.IwBernHk.bernHk},
\srcid{IwBernFaith.bernFaith},
\srcid{HkSandwich.hkSandwich},
\srcid{HkSum.hkSum},
\srcid{HkDualEx.hkDualEx}, and
\srcid{HkDualExp.hkDualExp}.
The passage from actual local group actions to these identities,
including cancellation of the positive integer index, is the proof
\srcid{PatchingModules.NilpotentLevelLowering.HkHypOf.locA} and
\srcid{hkHypOf}.
The maps and their two composites are constructed in
\srcid{MultiPhi.multiPhi}, \srcid{MultiPsi.multiPsi},
\srcid{MultiTwo.multiTwo}, and \srcid{TwoMul.twoMul}.
Their spherical inclusion and product trace in the global application
are written explicitly in
\srcid{ModularityArgument.ThresholdArguments.DenseThresholdCases.LevTwoOf.two\_gen}.
The arithmetic coefficient comparison (ON8) is the proof
\srcid{DenseThresholdCases.SymOfVieta.symOf\_alt}:
its specialization comparison is \srcid{symOf\_sym}, its coefficient
calculation is \srcid{symOf\_vieta}, and nonzero spherical specialization
is \srcid{symOf\_sph}.
All files are below
\url{https://github.com/ChristianKudeba/alpoge-rh-verification/tree/alpoge-source-7cbf7a0c/src/Development/AutomorphicArithmetic/AutomorphicVarieties}.

\subsection{The uniform auxiliary exponent and the finite-level module}
\label{sec:uniform-aux-exponent}

The integer exponent used before patching can be chosen to be
$n!$.  In rank three it is therefore $6$, chosen before the
host, the coefficient prime, the auxiliary level, and the
number of auxiliary places.

\begin{srlemma}[A polynomial vanishing on the torus support]
\label{lem:torus-uniform-nil}
Let a commutative group $T$ act by invertible commuting operators
on a complex vector space $W$ of dimension at most $d$.
Let $t_1,\ldots,t_r\in T$ and $F\in\mathbb C[X_1,\ldots,X_r]$.
If
\[
            F(\chi(t_1),\ldots,\chi(t_r))=0
\]
for every character $\chi$ occurring in a simultaneous
generalized-character space of $W$, then
$F(t_1,\ldots,t_r)^d=0$ on $W$.
The same exponent works on any direct sum of copies of this
$T$-module.
\end{srlemma}
\begin{proof}
A commuting family over an algebraically closed field has a
common eigenvector.  Apply this observation successively to
quotients to obtain a complete invariant flag.
The diagonal functions on $T$ are multiplicative characters,
since every group element is invertible.
The hypothesis says that the polynomial in the operators has
zero diagonal in this flag.  It is strictly upper triangular,
so its $d$-th power is zero.
On a direct sum the same identity holds on each summand.
\end{proof}

We state explicitly the local representation-theoretic input
used to take $d=n!$.
For a nonarchimedean local field $K$, a subgroup
$C\subseteq\kappa_K^\times$, and an irreducible smooth complex
representation $\sigma$ of $\operatorname{GL}_n(K)$, the space
$\sigma^{\operatorname{Iw}_C}$ is finite dimensional and
\[
                   \dim_\mathbb C
                      \sigma^{\operatorname{Iw}_C}\leq n!.
                                                               \tag{35}
\]
This is the exact bound \srcid{AlgebraTools.IwJqLine.jqLine},
using the printed geometric lemma and Bernstein--Zelevinsky
finiteness hypotheses.
Its mechanism is worth recording.
The canonical map from these invariants to the Borel Jacquet
module is injective: a vector in its kernel is killed by a
sufficiently positive Hecke operator, whereas the torus action
extending the positive Hecke action makes that operator
invertible.  A nonzero simultaneous torus weight, together
with Frobenius reciprocity, embeds $\sigma$ in a principal
series.  The Borel Jacquet module of this principal series has
a filtration with one one-dimensional quotient for each of
the $n!$ Bruhat cells, by the geometric lemma.
Exactness of the Jacquet functor in characteristic zero bounds
the Jacquet module of $\sigma$ by $n!$.
The injective canonical map gives (35).
If the invariant space is zero, (35) is immediate.
Thus this input has no dependence on the residue cardinality
or on $C$.

For an isotypic sum of copies of $\sigma$, invariants under the
compact group $\operatorname{Iw}_C$ are spanned by images of
$\sigma^{\operatorname{Iw}_C}$ under intertwining maps.
Indeed any finite expression can be averaged over this compact
group; smoothness reduces this averaging to a finite sum.
The torus action commutes with these intertwining maps.
Lemma~\ref{lem:torus-uniform-nil} and (35) therefore give the
same exponent $n!$ on an isotypic sum, regardless of its
multiplicity.  Restriction to invariants at other places
does not increase the required exponent.
This is the actual uniform-exponent argument in
\srcid{NilpotentLevelLowering.NilKill}.

\paragraph{The two auxiliary polynomial relations.}
At one auxiliary place write $Q$ for the residue cardinality
and $P,T$ for Frobenius and tame inertia matrices.
The characteristic-zero local data supplies a semisimplification
by characters $c_i$, and
\[
 P^{-1}TP=T^Q,\qquad
 \det(X-R(g))=\prod_i(X-c_i(g)).
\]
On the locus
\[
         \alpha_i=c_i(\mathrm{Frob}),\quad
         \alpha_i\ne\alpha_j,\quad
         \alpha_i\ne Q\alpha_j\quad(i\ne j),                    \tag{36}
\]
the following matrix identities hold:
\[
 \prod_i(P-\alpha_i)=0,\qquad
 \prod_{k\ne i}(P-\alpha_k)
                    (T-c_i(\tau))=0.                          \tag{37}
\]
Here the second product may be written in any order.

We include the matrix argument.
The character relation gives $c_i(\tau)^{Q-1}=1$; hence
$T^{Q-1}$ is unipotent.
Set $N=\log(T^{Q-1})$, using its finite nilpotent logarithm.
Then $P^{-1}NP=QN$.
The eigenvalues of $\operatorname{Ad}(P^{-1})$ are the ratios
$\alpha_j/\alpha_i$.  Condition (36), and $Q\ne1$, exclude $Q$,
so $N=0$.  Thus $T^{Q-1}=1$, and the tame relation gives $PT=TP$.
The eigenspaces of $P$ are one dimensional.
The character identities for the traces of $P^rT$,
$0\leq r<n$, and the invertible Vandermonde matrix in the
distinct $\alpha_i$, show that $T$ acts on the $i$-th
eigenspace by $c_i(\tau)$.  This proves (37).
This is the content of the finite matrix argument
\srcid{LocData.rel\_of\_chars}, also used by
\srcid{TLocDataT.locDataT\_alt} with the printed local--global
compatibility statement.

For all auxiliary places introduce commuting Hecke variables
$U_{ji}$ and $Y_{ji}$ and put
\[
 f_N=\prod_j\prod_i\prod_{i'\ne i}
            (U_{ji}-U_{ji'})(U_{ji}-Q_{Nj}U_{ji'}).
                                                               \tag{38}
\]
Let $r$ be either trace of a group element times the first
polynomial in (37), or trace of a group element times
\[
       \prod_{k\ne i}(P-U_{jk})(T-1-Y_{ji}).                    \tag{39}
\]
At every supported complex character of the relevant local
torus, the polynomial $f_Nr$ vanishes.
If $f_N$ vanishes this is immediate; otherwise its factors
give (36) and (37) gives the assertion.
The torus-support lemma proves
\[
                       (f_Nr)^{\,n!}M_N=0.                   \tag{40}
\]
All additional factors in (38), at the other auxiliary places,
commute with the operator at this place, so they do not
multiply the exponent by the number of places.

The passage back to the integral ordinary module uses the
actual separating family of classical automorphic
specializations: an element is zero if all its classical
images are zero.  Each specialization sends the auxiliary
operators to the normalized torus operators
\[
       U_{ji}\longmapsto Q_{Nj}^{\,i}\Theta_j(\varpi_i),
       \qquad 1+Y_{ji}\longmapsto\Theta_j(\delta_i).
\]
Their polynomial evaluations intertwine the integral and
complex actions.  Applying the already uniform identity
at every specialization and then separation proves (40)
integrally, with the same exponent $n!$.
This is \srcid{LevRelOf.relOf\_alt} and
\srcid{LocOkLevelModuleOfThetaData.level\_relation\_one\_two}.
In particular the exponent is fixed before $N$, before the
chosen relation, and before the element of $M_N$.

\begin{srlemma}[Projective norm and the augmentation variables]
\label{lem:projective-norm-control}
Let $(\Lambda,\mathfrak m)$ be a complete Noetherian local
ring with residue characteristic $\ell$ and
$\ell\in\mathfrak m$.  Let $\Delta$ be a finite $\ell$-group
and $M$ a finite projective $\Lambda[\Delta]$-module.
The norm
\[
       \mathcal N:M\longrightarrow M^\Delta,\qquad
                         x\longmapsto\sum_{\delta\in\Delta}\delta x
\]
is onto, with kernel the augmentation submodule
$\sum_{\delta\in\Delta}(\delta-1)M$.
If $\Delta$ is abelian, $\Lambda[\Delta]$ is local and
$M$ is finite free over it.
\end{srlemma}
\begin{proof}
For the regular module $\Lambda[\Delta]$, the invariant
elements are exactly the scalar multiples of
$\sum_\delta\delta$.  Its norm is multiplication of this
element by the augmentation.  Its kernel is therefore the
augmentation ideal, which is generated by $\delta-1$.
The same assertions hold on finite free modules.
A projective module is a direct summand of such a module;
norm, invariants, and augmentation submodules respect this
decomposition, proving the norm assertion.
For abelian $\Delta$, the reduction
$(\Lambda/\mathfrak m)[\Delta]$ is local:
its augmentation ideal is nilpotent and its quotient is the
residue field.  Thus the finite complete $\Lambda$-algebra
$\Lambda[\Delta]$ is local, with maximal ideal generated by
$\mathfrak m$ and $\delta-1$.
A finite projective module over a local ring is free.
\end{proof}

This proves the trace/control assertion at the auxiliary
level without dividing by $\#\Delta$, which is usually not a
unit.  For a chosen generating family $\delta_p$ one can use
only the operators $\delta_p-1$ in the kernel formula:
$\delta\delta'-1=\delta(\delta'-1)+(\delta-1)$.

\begin{srlemma}[Fixed truncations of the auxiliary group ring]
\label{lem:aux-trunc-free}
Let
\[
 \Delta_N=\prod_{j,i}\mathbb Z/\ell^{a_{Nj}}\mathbb Z,
                     \qquad a_{Nj}\geq N.
\]
Put $S=\Lambda[[Y_{ji}]]$, and let $M_N$ be finite projective
over $\Lambda[\Delta_N]$.  Identify $1+Y_{ji}$ with the
corresponding group generator.  Then for every fixed $k,c$,
and all sufficiently large $N$,
\[
 M_N/(Y^k+\mathfrak m^c)M_N
       \text{ is finite free over }
                      S/(Y^k+\mathfrak m^c).
                                                               \tag{41}
\]
Its rank may depend on $N$.
\end{srlemma}
\begin{proof}
There is an isomorphism
\[
 \Lambda[\Delta_N]\simeq
 S/\bigl((1+Y_{ji})^{\ell^{a_{Nj}}}-1\bigr)_{j,i}.
\]
The polynomials are monic and distinguished, so successive
power-series division proves the isomorphism as well as
finiteness and completeness.
For $1\leq r<k$,
\[
 v_\ell\binom{\ell^a}{r}\geq a-v_\ell(r),
\]
because $r\binom{\ell^a}{r}=\ell^a
\binom{\ell^a-1}{r-1}$.
Thus if $N\geq c+\max_{1\leq r<k}v_\ell(r)$, every displayed
relation belongs to $Y^k+\mathfrak m^c$.
Terms of degree at least $k$ already belong to $Y^k$.
The preceding lemma makes $M_N$ free over its group ring;
taking the specified quotient proves (41).
\end{proof}

\paragraph{Matching the common-module hypotheses.}
The module in (41) is the module entering the common
ultraproduct construction.  At $Y=0$, the projective norm
lemma gives its lower-level invariant module.
The two old/new-level maps, after the auxiliary root-splitting
algebra is introduced, have both composites equal to the
single discriminator $f_N$.  On the fixed base module these
maps use the same annihilator ideal $J$; their scalar
extensions do not introduce a new automorphic support ideal.
The two maps and their discriminator are the
\srcid{TwoMaps} input of \srcid{OpToTW.opCore}, assembled by
\srcid{LevTwoOf.twoOf\_alt}.

The source construction then takes the simultaneous
topologically nilpotent summand for $U_{ji}-u_{ji}$ and
$Y_{ji}$.  The complementary idempotent decomposition
commutes with all Hecke operators.  It preserves the
projective norm calculation and the identities (40).
Passing to the root-splitting algebra and completing sends
$Y$ to zero on the base module, and the two control maps
remain linear for this common specialization.
These are the operations \srcid{nilPart},
\srcid{freePart}, \srcid{quotPart}, and \srcid{spPart}
inside \srcid{NilpotentLevelLowering.OpToTW}.

For the fixed limit auxiliary elements, their finite-level
trace relations converge to the relations defining the
common ideal $I$.
Also $Q_{Nj}\equiv1\pmod{\ell^N}$, so (38) converges to the
fixed discriminant $f$.
Consequently (40), with the fixed exponent $e=n!$, says
exactly: for every generator $r$ of $I$ and every precision
$c$, for all sufficiently large $N$, $(fr)^e$ is congruent
modulo $\mathfrak m_R^c$ to an element annihilating $M_N$.
This is the uniform approximate annihilator condition of
the common patched-module theorem.

The exact assembly in \srcid{THasAutOf} therefore chooses
\srcid{TWMaps} with exponent \srcid{n.factorial}.
For rank three, $e=6$ throughout.  The free truncations
(41), the norm/control statement, and the uniform exponent
(40) have all been proved here.
The ordinary function-module construction, projectivity and old/new
maps are supplied earlier in Section~\ref{sec:sr-ordinary-control}.
For the actual arithmetic specialization one must still construct
the attached Galois determinant and verify its classical separation,
local--global compatibility, and the same universal ring and limit
data. These requirements are not consequences of the existence of
an abstract deformation ring.

\subsection{The cohomological passage to the Taylor--Wiles limit}
\label{sec:coh-at-tw-limit}

This section proves the limit conversion \srcid{CohLim}.
It is separate from the arithmetic estimates producing the
finite-level hypotheses \srcid{B1} and \srcid{B2}.

Keep $A_0=k[[T]]$, $F=k((T))$, $q=\#k$, the paired point $\rho$,
and the chosen odd classes $\kappa_1,\ldots,\kappa_g$ on the
\emph{original} group $\Gamma_0$.
Their original-group generating property is an explicit premise:
for every integral odd Selmer cocycle $z_0$ on $\Gamma_0$ there
are $s\geq0$ and $a_j\in A_0$ such that
$T^sz_0-\sum_j a_j\kappa_j$ is an integral odd coboundary.
This is exactly \srcid{Kap.gen}; independence is
\srcid{Kap.indep}.  It is not a spanning assumption on
$\Gamma_\infty$, which is what the proof establishes.
For $m\geq1$ write $A_m=A_0/T^m$.
Use the even, odd, first-order and second-order Selmer conditions
specified above, including their local correcting matrices.
Let $\Gamma_N\to\Gamma_0$ be the auxiliary groups and
$\Gamma_\infty\to\Gamma_0$ the group denoted by
$\Lambda.G$ in \srcid{Lim}.
The limit data include the following facts; these are the
specific convergence properties used in the proof.
\begin{enumerate}
\item A topologically finitely generated profinite group $P$
has continuous surjections $\pi_N:P\to\Gamma_N$ and
$\pi_\infty:P\to\Gamma_\infty$.
They preserve the sign characters and a chosen complex
conjugation.  The index-two sign kernel $P^+$ is therefore
also topologically finitely generated.
\item For every open normal $V\subseteq P$, eventually along
the chosen nonprincipal ultrafilter,
\[
       (\ker\pi_N)V=(\ker\pi_\infty)V.
                                                               \tag{30}
\]
The maps $f_N:P\to\Gamma_0$ and $f_\infty:P\to\Gamma_0$
obtained by composition agree modulo every fixed finite
quotient of $\Gamma_0$, eventually and uniformly on $P$.
\item The local maps, the prescribed unramified local maps,
and the auxiliary elements $\phi_{Nj},\tau_{Nj}$ converge
in the same finite-quotient sense on $P$.
More precisely their values have lifts to $P$ which agree
modulo every prescribed open normal $V$ eventually,
uniformly over the local group.  This is the relation
\srcid{Close} and the assertions
$\Lambda.\mathrm{phi\_lim}$,
$\Lambda.\mathrm{phiu\_lim}$,
$\Lambda.\mathrm{ph\_lim}$,
$\Lambda.\mathrm{ta\_lim}$.
\item After the \srcid{ProL} quotient, each
$\ker(\Gamma_N\to\Gamma_0)$ is contained in the closed normal
subgroup generated by the $\tau_{Nj}$.
At the limit
\[
 h_\infty(\tau_j)=1,\quad h_\infty(\phi_j)=s_j,\quad
                         \phi_j\tau_j=\tau_j\phi_j.             \tag{31}
\]
The matrix $\rho(s_j)$ has $n$ distinct eigenvalues in $F$.
\end{enumerate}
All these are statements about the existing limit object.
For example (30) is $\Lambda.\mathrm{ker\_lim}$, and (31)
is $\Lambda.\mathrm{p\_ta}$,
$\Lambda.\mathrm{p\_ph}$ and $\Lambda.\mathrm{comm}$.

The exact finite-level inputs, with constants chosen before $m,N$,
are these:
\[
 \begin{split}
 \#\operatorname{Sel}_{1,N}(A_m)
 &\leq
 \#\operatorname{Bd}_{1,N}(A_m)\,q^{\,mD+C_1},\\
 &\hspace{5mm}\operatorname{Sel}_{1,N}(A_m)\text{ is finite},
 \end{split}                                                    \tag{32}
\]
eventually in $N$ for every fixed $m$; and, for every $g+1$
tuples of even Selmer cocycles $z_{ij}$ over $A_m$,
\[
 \begin{gathered}
 \exists\,c_1,\ldots,c_{g+1}\in A_0
       \text{ with at least one }c_i\in A_0^\times,\quad
 \exists\,\Xi\in\operatorname{Sel}_{2,N}(A_m),\\
       \varepsilon_j'(\Xi)=T^{C_2}\sum_i c_i z_{ij}
                             \quad(1\leq j\leq g).
 \end{gathered}                                                   \tag{33}
\]
For each fixed $m$, (33) holds for all such tuples at all
sufficiently large $N$.
These are \srcid{TWGrp.B1} and \srcid{TWGrp.B2}.
The uniform exponent $C_2$ and the primitive coefficient vector
in (33) are essential.

\begin{srproposition}[The three limit bounds]
\label{prop:tw-limit-coh}
Under these hypotheses the integral first-order Selmer rank
is at most $D$; every $g+1$ tuples of integral even cocycles
have a nonzero $A_0$-linear combination lifting to an integral
second-order cocycle; and every integral odd Selmer cocycle
on $\Gamma_\infty$ is an $A_0$-linear combination of the
$\kappa_j$ modulo a coboundary after multiplication by some
power of $T$.
These are exactly $K1(h_\infty,\kappa,D)$,
$K2(h_\infty,\kappa)$ and $K3(h_\infty,\kappa)$.
They imply the three Laurent hypotheses of
Theorem~\ref{thm:selmer-paired-pres}.
\end{srproposition}

\begin{proof}
\emph{Finite coefficient Selmer spaces transfer to the limit.}
For a fixed finite coefficient module and a fixed continuous
action, the continuous cocycles on a topologically finitely
generated profinite group form a finite set.
Indeed a cocycle is determined by its values on finitely many
topological generators: its values on the abstract subgroup
they generate follow from the cocycle identity, and continuity
gives the rest.
Equivalently use the homomorphism into the finite semidirect
product of the module and its action group.

A locally constant function $c:P^+\to Z$ factors through a
surjection $\pi:P^+\to H$ precisely when $\ker\pi$ lies in its
right stabilizer
\[
             \operatorname{Stab}(c)
             =\{u:c(xu)=c(x)\ \forall x\}.
\]
This stabilizer is open, since $P^+$ is compact.
Its normal core, viewed in $P$, is open normal.
Thus (30) says that $c$ factors through $\pi_N$ precisely
when it factors through $\pi_\infty$, eventually.
There are only finitely many cocycles on $P^+$ for the chosen
finite action, so one may impose this assertion for all of
them at once.
The actions themselves agree for large $N$ by the finite
quotient convergence of $f_N$.
The local convergence property in item 3 makes the descended
functions agree on all specified local and unramified groups.
The same correcting matrix therefore verifies the ordinary
condition on both groups.
Polarization is preserved by the fixed conjugation.
Consequently the finite coefficient Selmer sets on
$\Gamma_N$ and $\Gamma_\infty$ are identified by their pullbacks
to $P^+$, for sufficiently large $N$.
The coboundary sets are identified as well.
This proves (32)--(33) on $\Gamma_\infty$, with the same
constants $C_1,C_2$.

\emph{A fixed loss of precision for independent classes.}
Let $z_1,\ldots,z_e$ be integral first-order cocycles independent
over $A_0$ modulo integral first-order coboundaries.
There is $C\geq0$ such that
\[
 \sum_i a_i z_i\text{ is a coboundary modulo }T^m
       \quad\Longrightarrow\quad
                         a_i\in T^{m-C}A_0\quad(m\geq C).       \tag{34}
\]
Here is a finite-module proof of this assertion.
Let $V$ be the finite free $A_0$-module of allowed
skew-transpose first-order correcting matrices.
On $M=A_0^e\oplus V$ consider the linear operators
\[
        L_\gamma(a,Y)=\sum_i a_i z_i(\gamma)-dY(\gamma).
\]
Their span in the finite module
$\operatorname{Hom}_{A_0}(M,\operatorname{Mat}_n(\mathbb D_g(A_0)))$
is finitely generated.  Choose finite generators and assemble
them into one linear map $L:M\to M'$ with $M'$ finite free.
The common kernel is the kernel of all $L_\gamma$.
Independence of the classes says that each coordinate
functional $a\mapsto a_i$ vanishes on this kernel.
Smith normal form over the DVR $A_0$ supplies $C$ with
\[
       L(x)\in T^mM'
                 \ \Longrightarrow\
       x\in\ker L+T^{m-C}M.
\]
A coboundary matrix modulo $T^m$ lifts to $V$:
lift its upper-triangular entries and reflect them with the
required minus sign, preserving the first-order parity.
Applying the last implication to this lift proves (34).
This is the \srcid{CoefficientsBoundedOfSelmerBoundary}
argument, expressed over a DVR instead of its more general
Artin--Rees formulation.

Choose representatives in $A_0$ for $A_0/T^{m-C}$.
Their $e$-tuples give distinct classes modulo the first-order
coboundaries by (34).  Hence
\[
        q^{(m-C)e}
       \leq
       \frac{\#\operatorname{Sel}_{1,\infty}(A_m)}
            {\#\operatorname{Bd}_{1,\infty}(A_m)}
       \leq q^{\,mD+C_1}.
\]
Letting the integer $m$ grow proves $e\leq D$.  This is $K1$.

\emph{Second-order lifting survives with its local conditions.}
For integral tuples $z_{ij}$, (33) gives primitive coefficient
vectors $c_m$ and second-order Selmer cocycles $\Xi_m$ modulo
$T^m$.  The ring $A_0$ is compact.
Moreover the set of second-order Selmer cocycles modulo
$T^j$ is finite for every fixed $j$, by the finite-cocycle
argument above and finite generation of $\Gamma_\infty$.

Take an ultralimit of $c_m$ and of all the reductions
$\Xi_m\bmod T^j$.  The latter reductions stabilize as whole
functions, not merely pointwise, because they belong to a
finite set of continuous cocycles.
They are compatible and define an integral continuous
cocycle $\Xi$.
Its coefficient, cocycle, polarization and unramified identities
hold at every precision and hence integrally.

The local correcting matrices also survive.
At each local place, lift the correcting matrices for
$\Xi_m$ arbitrarily to the finite free $A_0$-module of all
matrices and take their ultralimit in this compact module.
Parity and the equations for zero lower entries and the
specified zero diagonal entries are closed conditions.
Thus the limit is an allowed correcting matrix for $\Xi$.
This proves the ordinary condition without discarding it
when passing to the limit.

Pass to an ultrafilter set on which one fixed coordinate of
$c_m$ is a unit; one exists because there are only $g+1$ coordinates.
The limiting vector $c$ still has a unit coordinate.
The defining equations become
\[
                    \varepsilon_j'(\Xi)
                        =T^{C_2}\sum_i c_i z_{ij}.
\]
The vector $T^{C_2}c$ is nonzero since $A_0$ is a domain.
This proves $K2$.

\emph{The kernel is generated by the auxiliary inertia.}
Let $K$ be the closed normal subgroup of $\Gamma_\infty$
generated by all $\tau_j$.
We claim $\ker h_\infty\subseteq K$.
It suffices to test a finite quotient
$Q=\Gamma_\infty/W$ with $K\subseteq W$ open normal.
Let $V=\pi_\infty^{-1}(W)$.
Equation (30) makes $\ker\pi_N\subseteq V$ eventually, giving
$\Gamma_N\to Q$ through the common cover.
The convergence of the $\tau_{Nj}$ makes their images in $Q$
trivial.
Item 4 therefore makes this homomorphism descend to a
homomorphism $\theta_N:\Gamma_0\to Q$.

There are only finitely many continuous homomorphisms
$\Gamma_0\to Q$, because $\Gamma_0$ is a quotient of the
topologically finitely generated group $P$.
Take the intersection $U$ of their kernels.
Finite quotient convergence of $f_N$ to $f_\infty$ makes them
agree modulo $U$ eventually, so
$\theta_Nf_N=\theta_Nf_\infty$.
But the first of these is precisely
$P\to\Gamma_\infty\to Q$.
For $x\in\ker h_\infty$, choose a lift in $P$; its image under
$f_\infty$ is $1$, and therefore its image in $Q$ is $1$.
This proves the claim.  It is the substance of
\srcid{KerOo.ker\_oo}.

\emph{Odd classes descend.}
Let $z$ be an integral odd Selmer cocycle on $\Gamma_\infty$.
The cocycle identities for $\phi_j\tau_j=\tau_j\phi_j$, and
$h_\infty(\tau_j)=1$, imply
\[
                       \rho(s_j)z(\tau_j)
                                  =z(\tau_j)\rho(s_j).
\]
The matrix $\rho(s_j)$ is even and has distinct eigenvalues
over $F$, while $z(\tau_j)$ is odd.
The centralizer of a matrix with distinct eigenvalues is
commutative.  Its even block idempotent commutes with
$\rho(s_j)$, hence also with $z(\tau_j)$.
An odd matrix commuting with this block idempotent is zero.
Thus $z(\tau_j)=0$.

The subgroup on which both $h_\infty$ and $z$ vanish is closed
and normal.  Normality follows from the cocycle identity;
its stability under complex conjugation follows from
$z(\overline g)=-z(g)^{\mathsf t}$.
The kernel-generation claim therefore shows that $z$ vanishes
on $\ker h_\infty$, and $z$ descends to $\Gamma_0$.
Its continuity and all Selmer conditions descend:
the group map is a continuous quotient, the local maps
commute with it, and the same local correcting matrices apply.
The generating property of the original $\kappa_j$ now gives
$T^s z-\sum_j a_j\kappa_j$ a coboundary for some $s\geq0$ and
$a_j\in A_0$.  This proves $K3$.

Finally these integral conclusions imply the three Laurent
inputs.  For a finite family of bounded Laurent cocycles,
clear one common power of $T$ in the functions and in their
local correcting matrices; the distinguished local set is
finite in the arithmetic application.
This gives integral Selmer cocycles.
Apply $K1$ to an independent integral family, $K2$ to the
scaled even tuples, and $K3$ to a scaled odd cocycle.
After inverting $T$, their conclusions are exactly the
Laurent rank bound, nonzero second-order combination, and
spanning by $\kappa_j$.
No dimension or local condition is changed by this scalar
normalization.
\end{proof}

\paragraph{Sources and the remaining arithmetic input.}
Finite coefficient transfer is
\srcid{LcOfAdic.t1}, with
\srcid{EventuallyFactorThroughTowerIffLimit} and
\srcid{EventuallyFactorizationCompatibleWithLocalMaps}.
The first-order argument is
\srcid{FirstCohomologyBoundOfBoundaryBound}.
The second-order argument is
\srcid{SecondCohomologyBoundOfBoundaryBound},
\srcid{SelmerSecondExistsUltrafilterLimit}, and
\srcid{SelmerSecondClosedUnderLimits}.
Odd descent is \srcid{KoszulThreeOfDet}, \srcid{OdComm},
\srcid{CocKer}, and \srcid{KerOo}.
The assembly is \srcid{AutomorphicDetection.CohLim.cohLim}.
The modules are under
\srcid{Development.AutomorphicArithmetic.AutomorphicVarieties.%
PatchingModules} at the frozen source tag.
Thus $K1,K2,K3$ require no additional theorem once the actual
limit object, \srcid{ProL}, distinct auxiliary eigenvalues,
and the uniform finite-level bounds (32)--(33) have been
produced.  Their arithmetic production is not asserted by
the present limit argument.
